\newif\ifuseseminar
\useseminartrue
\input{../../latex_common/header.tex.frag}

\title{Temas para Seminários de Física Matemática II}

\begin{document}

\section*{Seminário 1: Ortogonalidade e Normalização dos Harmônicos Esféricos}

\textbf{Objetivo:} Discutir a estrutura matemática dos harmônicos esféricos focando na separação angular, ortogonalidade e normalização.

\begin{itemize}
	\item Separação da parte angular da equação de Laplace;
	\item Solução para a variável \( \phi \): série de Fourier com base ortonormal \( e^{im\phi} \);
	\item Equação para \( \theta \): surgimento dos polinômios associados de Legendre \( P_\ell^m(\cos\theta) \);
	\item Ortogonalidade dos \( P_\ell^m \) e dos harmônicos esféricos \( Y_\ell^m(\theta,\phi) \);
	\item Normalização e papel dos harmônicos como base ortonormal em \( L^2(S^2) \).
\end{itemize}

\vspace{1cm}

\section*{Seminário 2: Separação de Variáveis em Coordenadas Cartesianas e Esféricas}

\textbf{Objetivo:} Apresentar o método de separação de variáveis em duas geometrias, com foco na ortonormalidade das soluções em coordenadas cartesianas.

\begin{itemize}
	\item Método de separação de variáveis aplicado à equação de Laplace;
	\item Coordenadas cartesianas: resolução no interior de um cubo com condições de Dirichlet homogêneas;
	\item Construção de uma base ortonormal de funções tipo seno;
	\item Propriedades de ortogonalidade e normalização (produto interno em \( L^2 \));
	\item Coordenadas esféricas: separação formal da equação (sem resolver);
	\item Interpretação das equações resultantes para \( r \), \( \theta \), \( \phi \).
\end{itemize}

\vspace{1cm}

\section*{Seminário 3: Equação de Poisson, Função de Green e Teorema da Unicidade}

\textbf{Objetivo:} Apresentar a solução formal da equação de Poisson usando funções de Green e discutir o teorema da unicidade.

\begin{itemize}
	\item Equação de Poisson: \( \nabla^2 \Phi = -\rho/\epsilon_0 \);
	\item Solução formal com função de Green;
	\item Interpretação física da função de Green (resposta à fonte pontual);
	\item Exemplo: função de Green no espaço livre;
	\item Declaração e implicações do teorema da unicidade (Dirichlet e Neumann);
	\item Consequências físicas: garantia de que a solução obtida por métodos formais é única.
\end{itemize}

\vspace{1cm}

\section*{Seminário 4: Parte Radial da Equação de Schrödinger para o Hidrogênio}

\textbf{Objetivo:} Analisar a equação radial do átomo de hidrogênio, com ênfase nos aspectos matemáticos que levam à quantização da energia.

\begin{itemize}
	\item Separação de variáveis na equação de Schrödinger com potencial eletrostático;
	\item Forma da equação radial efetiva com o termo centrífugo;
	\item Análise assintótica para \( r \to 0 \) e \( r \to \infty \);
	\item Mudança de variáveis e definição da função \( u(r) = r R(r) \);
	\item Método de Frobenius e imposição de que a série termine (polinômio de Laguerre associado);
	\item Surgimento da quantização da energia \( E_n \sim -1/n^2 \).
\end{itemize}

\end{document}
